\SourceLeaf{196}%
% A544-P0196-D01
Auch für die vierte, fünfte und sechste Million hat der Verf. die Materialien bereits grösstentheils vorräthig, und er erbietet sich, diese Fortsetzung zu liefern, wenn der Verleger durch einen hinreichenden Absatz der drei ersten Millionen aufgemuntert wird. Es wäre in der That sehr zu beklagen, wenn die Früchte einer so mühsamen und nützlichen Arbeit der Welt entzogen werden sollten.

% A544-P0196-D02
\Needspace{8\baselineskip}
\pdfbookmark[1]{Erchinger, Siebenzehneck. Anzeige 1825 December 19}{notice-1825-12-19}
\begin{center}
\rule{.86\linewidth}{.3pt}\par
Göttingische gelehrte Anzeigen.\quad 1825 December 19.\par
\rule{.86\linewidth}{.3pt}
\end{center}
Der Königl. Societät ist abseiten des Herrn \textsc{Erchinger} zu Thuningen im Königreich Würtemberg eine kleine Abhandlung vorgelegt worden, welche die
\begin{center}\emph{Geometrische Construction des regelmässigen Siebenzehnecks}\end{center}
zum Gegenstande hat. Die Allgemeine Theorie der regelmässigen Vielecke hat bekanntlich durch die innige Verbindung, in welche sie mit der höhern Arithmetik gebracht ist, eine neue Gestalt und Erweiterung erhalten; ein, wenn gleich verhältnissmässig nur kleiner Theil derselben ist die Theorie derjenigen Vielecke, die sich geometrisch beschreiben lassen. Seit dem Zeitalter der Griechen wusste man, dass das Dreieck, Fünfeck, Funfzehneck und alle diejenigen Vielecke, welche durch Verdopplung oder wiederholte Verdopplung der Seitenzahl aus diesen entspringen, jene Eigenschaft haben, und man glaubte, behauptete auch wohl ausdrücklich, dass dieses die einzigen seien. Die höhere Arithmetik hat gelehrt, dass dieses ein Irrthum war: indem sie die wahren Quellen der ganz allgemeinen Theorie offen legte, ergab sich von selbst, dass es ausser den genannten Vielecken noch unzählige andere gibt, die geometrisch construirt werden können, von denen das Siebenzehneck das einfachste ist. Die Ueberlegenheit der Analyse, welche das Allgemeinste, wie das Besondere mit gleicher Leichtigkeit umfasst, über die Geometrie, die immer beim Besondern stehen bleiben muss, beim Fortschreiten von den einfachern Fällen zu den zusammengesetztern durch stets vergrösserte Verwicklung aufgehalten wird, und jenen den bekannten nächsten Fall schwerlich jemals ohne fremde Hülfe erreicht hätte, zeigt sich dabei im hellsten Lichte. Inzwischen ist es immer wichtig, interessant und wünschenswerth, dass auch die rein geometrischen Behandlungen fortwährend cultivirt werden, und dass die 
\SourceLeaf{197}%
% A544-P0197-D01
Geometrie wenigstens einen Theil der neuen Felder, die die Analyse erobert, sich aneigne. Ref. ist nicht bekannt, dass bisher jemand die Construction des Siebenzehnecks öffentlich behandelt hätte, ausser Herrn \textsc{Pauker} in den Schriften der Kurländischen Gesellschaft und in seiner Geometrie. Verschieden davon und mehr im rein geometrischen Geiste durchgeführt ist die von Hrn \textsc{Erchinger}, welche in Folgendem besteht. (Die dazu gehörige Figur, eine gerade Linie, auf welcher der Folge nach die Punkte $DBGAIFCE$ liegen, kann jeder sich selbst zeichnen.) Eine nach Gefallen angenommene gerade Linie $AB$ verlängere man rückwärts und vorwärts nach $C$ und $D$ so, dass $AC\times BC=AD\times BD=4AB\times AB$ werden; ferner bestimme man die Punkte $E$, $G$ an beiden Seiten der verlängerten Linie $CA$ so, dass $AE\times EC=AG\times CG=AB\times AB$, und den Punkt $F$ auf der Seite $A$ der verlängerten Linie $BA$ so, dass $AF\times DF=AB\times AB$ wird; endlich theile man $AE$ in $I$ so, dass $AI\times EI=AB\times AF$ werde, wo $AI$ der kleinere, und $EI$ der grössere Abschnitt von $AE$ ist. Man mache dann ein Dreieck, in welchem zwei Seiten jede $=AB$, die dritte $=AI$ wird. Beschreibt man um dieses Dreieck einen Kreis, so wird $AI$ die Seite des in den Kreis beschriebenen regelmässigen Siebenzehnecks sein.

% A544-P0197-D02
Wenn man die Richtigkeit dieser Construction durch die Vergleichung mit der in den \emph{Disquisitiones Arithmeticae} Art. 354 als ein Beispiel aufgestellter Theorie des Siebenzehnecks prüft, so bemerkt man leicht, dass jene nichts anders ist, als die geometrische Uebersetzung derjenigen Gleichungen, auf welche die Anwendung der allgemeinen Theorie führt: in der That sind die Entfernungen der Punkte $C$, $D$, $E$, $F$, $G$, $I$ von $A$ nichts anderes, als die Grössen, die a. a. O. mit $(8\pdot1)$, $(8\pdot3)$, $(4\pdot1)$, $(4\pdot3)$, $(4\pdot9)$, $(2\pdot1)$ bezeichnet sind, wenn man das positive und negative Zeichen durch die Lage ausdrückt, und die Entfernung des Punktes $B$ von $A$ in eben dem Sinn genommen $=-1$ setzt. Allein das eigentlich Verdienstliche der Abhandlung des Hrn. \textsc{Erchinger} besteht nicht sowohl in der Aufstellung der Construction selbst, da die Analyse bereits den einfachsten Weg vorgezeichnet hatte, als in der rein geometrischen Begründung ihrer Richtigkeit, und diese ist mit so musterhafter mühsamer Sorgfalt, alles nicht rein Elementarische zu vermeiden, durchgeführt, dass sie dem Verf. zur Ehre gereicht, und den Wunsch veranlasst, dass sein in der That nicht gemeines mathematisches Talent alle Aufmunterung finden möge.

\SourceLeaf{198}%
% A544-P0198-D01
\Needspace{8\baselineskip}
\pdfbookmark[1]{Seeber, positive ternäre quadratische Formen. Anzeige 1831 Juli 9}{notice-1831-07-09}
\begin{center}
\rule{.86\linewidth}{.3pt}\par
Göttingische gelehrte Anzeigen.\quad 1831 Juli 9.\par
\rule{.86\linewidth}{.3pt}
\end{center}
\emph{Untersuchungen über die Eigenschaften der positiven ternären quadratischen Formen von} \textsc{Ludwig August Seeber}, \emph{Dr. der Philosophie, ordentl. Professor der Physik an der Universität in Freiburg. Freiburg im Breisgau} 1831. (248 S. in 4.)

% A544-P0198-D02
Die Functionen zweier unbestimmten Grössen $x$ und $y$ von der Gestalt $axx+2bxy+cyy$, wo $a,b,c$ bestimmte ganze Zahlen vorstellen, bilden bekanntlich unter dem Namen der \emph{quadratischen Formen}, oder, wo eine weitere Unterscheidung erforderlich wird, der \emph{binären quadratischen Formen}, einen der interessantesten und reichhaltigsten Gegenstände der höheren Arithmetik. Die dabei zunächst vorkommenden Aufgaben: zu entscheiden, ob eine solche gegebene Form eine andere $a'x'x'+2b'x'y'+c'y'y'$ unter sich begreift, d. i. durch eine Substitution $x=\alpha x'+\beta y'$, $y=\gamma x'+\delta y'$, in welcher $\alpha,\beta,\gamma,\delta$ ganze Zahlen sind, in dieselbe verwandelt werden kann; ob eine solche Relation zweier Formen eine gegenseitige ist, wo die Formen äquivalent heissen; ferner in beiden Fällen alle möglichen Umformungen der einen in die andere anzugeben; endlich alle möglichen Darstellungen einer gegebenen ganzen Zahl durch eine gegebene Form vermöge ganzer Werthe der unbestimmten Grössen aufzufinden --- diese Aufgaben sind in den \emph{Disquisitiones Arithmeticae} vollständig aufgelöset, machen aber von dem die quadratischen Formen betreffenden Abschnitte dieses Werks nur den bei weitem kleineren Theil aus. Die darauf folgenden feineren Untersuchungen erfordern\textsuperscript{.} zum Theil eine vorläufige Bearbeitung eines um eine Stufe höheren und viel grössere Schwierigkeiten darbietenden Feldes, nemlich der Lehre von ähnlichen Functionen dreier unbestimmter Grössen $x,y,z$, welche also die Gestalt haben $axx+byy+czz+2a'yz+2b'xz+2c'xy$, und ternäre quadratische Formen heissen. Die Auflösung der diese ternären Formen betreffenden Hauptaufgaben ist in dem erwähnten Werke entwickelt, jedoch nur so weit, als zu dem angezeigten Zwecke nothwendig war. Nach einem Zwischenraum von dreissig Jahren hat nun der Verfasser des vorliegenden Werks zuerst diese Untersuchungen wieder aufgenommen, und in Beziehung auf die eine Hauptgattung der ternären Formen, nemlich die positiven, dasjenige was in den \emph{Disquisitiones }
\SourceLeaf{199}%
% A544-P0199-D01
\emph{Arithmeticae} unvollendet gelassen war, zur Vollständigkeit gebracht. Für diejenigen, welche aus der höheren Arithmetik ein tieferes Studium gemacht haben, würden wir dasjenige, was in dem vorliegenden Werke Neues geleistet ist, mit wenigen Worten bezeichnen können; allein, um auch andern verständlich zu sein, müssen wir uns etwas mehr Ausführlichkeit verstatten, und wir thun dies um so lieber, da diese Untersuchungen auch ausserhalb des Gebietes der höheren Arithmetik ein eigenthümliches Interesse haben.

% A544-P0199-D02
Die Eigenschaften einer binären Form $axx+2bxy+cyy$ hängen vornehmlich von der Zahl $bb-ac$ ab, welche daher der Determinant jener Form heisst. Zwei äquivalente Formen haben allemal gleiche Determinanten. Allein nicht alle Formen, die einen gegebenen Determinanten haben, sind darum schon äquivalent, vielmehr zerfallen solche Formen in eine kleinere oder grössere, aber stets endliche Anzahl von Classen, so dass die zu einerlei Classe gehörigen unter sich äquivalent, die zu verschiedenen Classen gehörenden hingegen nicht äquivalent sind. Durch Formen, deren Determinant positiv ist, lassen sich ohne Unterschied positive und negative Zahlen darstellen; hingegen durch Formen mit negativem Determinanten sind nur solche Zahlen darstellbar, welche mit $a$ und $c$ einerlei Zeichen haben, daher hier positive und negative Formen unterschieden werden. Die einfachsten Formen in jeder Classe haben bestimmte Kriterien, heissen reducirte Formen, und können als Repräsentanten der ganzen Classe betrachtet werden.

% A544-P0199-D03
Aehnliche Verhältnisse in Beziehung auf die ternären Formen sind in den \emph{Disquisitiones Arithmeticae} nachgewiesen. Determinant der ternären Form
\[
axx+byy+czz+2a'yz+2b'xz+2c'xy
\]
heisst die Zahl
\[
aa'a'+bb'b'+cc'c'-abc-2a'b'c'
\]
Auch hier ist zur Aequivalenz zweier Formen die Gleichheit der Determinanten erforderlich, aber nicht zureichend, sondern sämmtliche Formen mit einem bestimmten Determinanten zerfallen in eine endliche Anzahl von Classen, in deren jeder die einfachsten Formen reducirte heissen können und alle übrigen gleichsam repräsentiren. Mit dem Unterschiede zwischen positiven und negativen Formen verhält es sich aber hier anders, als bei den binären Formen. Für jeden gegebenen Determinanten, er sei positiv oder negativ, gibt es theils Formen, durch welche
\SourceLeaf{200}%
% A544-P0200-D01
ohne Unterschied positive und negative Zahlen darstellbar sind (indifferente Formen), theils solche Formen, durch die entweder nur positive oder nur negative Zahlen sich darstellen lassen (positive oder negative Formen); allein positive Formen gibt es nur für negative Determinanten, und negative nur für positive. Uebrigens ist es von selbst klar, dass die Qualification einer Form, insofern sie indifferent, positiv oder negativ ist, zugleich der ganzen Classe, zu welcher sie gehört, zukommt. Das vorliegende Werk beschränkt sich auf die positiven Formen, deren Determinanten also negativ sein müssen: offenbar findet aber alles, was von diesen gilt, von selbst seine Uebertragung auf die negativen Formen, während die in dem Werke ganz ausgeschlossenen indifferenten Formen eine ganz abweichende Behandlung erfordern.

% A544-P0200-D02
In den \emph{Disquisitiones Arithmeticae} war, wie schon erwähnt ist, die Theorie der ternären Formen nur so weit entwickelt, als für den dortigen Zweck nöthig war, und daher die Aufgabe, die Aequivalenz zweier gegebenen ternären Formen zu entscheiden, noch nicht in vollständiger Allgemeinheit aufgelöset. Zwar war daselbst gezeigt, wie man zu jeder vorgegebenen Form eine äquivalente der einfachsten Art finden, und dass es solcher reducirten Formen für jeden gegebenen Determinanten nur eine endliche Anzahl geben könne; allein da es in jeder Classe mehrere solcher reducirten Formen gibt, die sich nicht in allen Fällen \emph{sogleich} als äquivalent ergeben, so fehlte noch ein Kriterium, woran man die Aequivalenz oder Nicht-Aequivalenz solcher Formen mit Gewissheit erkennen kann. Dieses Bedürfniss hat nun der Verfasser des vorliegenden Werks in Beziehung auf die positiven Formen vollständig und mit musterhafter Gründlichkeit gehoben. Sein Verfahren ist übrigens etwas anders eingekleidet, als wir die Sache so eben ausgesprochen haben, und wie sie sich verhalten müsste, wenn man in den Begriff der reducirten positiven Formen nur die wesentlichsten Bedingungen der grössten Einfachheit aufnimmt, welche in dem Fall der positiven Formen die sind, dass die (ihrer Natur nach positiven) Zahlen $a,b,c$ nicht kleiner sein dürfen, als respective $b'$ oder $c'$, $a'$ oder $c'$, $a'$ oder $b'$ ohne Rücksicht auf die Zeichen. Herr \textsc{Seeber} hat nemlich dem Begriff der reducirten Formen noch solche Modificationen hinzugesetzt, dass es in jeder Classe immer nur Eine der Art geben kann, Eine aber geben muss. Wegen eines schönen von Herrn \textsc{Seeber} durch Induction gefundenen weiter unten noch zu erwähnenden Theorems führen wir hier die Hauptbedingungen, welche Hr. S. in den Begriff der reducirten Formen 
